\section{Quantitative consequences of the compatible rectangles}
\label{f:sec:quantitative-consequences}

\subsection[Completion in the global \texorpdfstring{\(H^1\)}{H1} mild class]
{Completion in the global \(H^1\) mild class}

The forced construction begins with velocity data in
\(H^\infty_\sigma(\R^3)\), not only with Schwartz data, while retaining the
one prescribed force throughout. The datum-generated full-interval
intersection, its finite rectangle projections, and same-force continuation
transfer, after positive-time smoothing, to the strong \(H^1\) class for that
force.

\begin{theorem}[Global \(H^1\) completion]
\label{f:thm:qc-global-H1-completion}
Let \(\nu>0\), let \(a\in H^1_\sigma(\R^3)\), and prescribe
\(f\in C^\infty([0,\infty);\mathcal S(\R^3;\R^3))\). Then there is a unique
pressure-normalized mild solution
\begin{equation}
 u^a\in C([0,\infty);H^1_\sigma)
 \cap L^2_{\rm loc}([0,\infty);H^2_\sigma),
 \qquad
 u^a\in C^\infty((0,\infty)\times\R^3),
 \label{f:eq:qc-global-H1-class}
\end{equation}
of the forced three-dimensional Navier--Stokes equations with datum \(a\),
viscosity \(\nu\), and the same prescribed force \(f\).  Its pressure is
\[
 p^a=R_iR_j(u_i^a u_j^a)-(-\Delta)^{-1}\nabla\cdot f
\]
in distributions.  The solution obeys the forced kinetic-energy identity on
every finite interval.
\end{theorem}

\begin{proof}
Set \(g:=\Pp f\), and write
\[
 G_0:=\norm g_{L^2(0,1;L^2)},
 \qquad
 R_{a,f}:=1+\norm{\nabla a}_2^2+2\nu^{-1}G_0^2.
\]
Write
\[
 Y(t):=\norm{\nabla u(t)}_2^2,
 \qquad
 Z(t):=\norm{\Delta u(t)}_2^2,
 \qquad
 c_1:=\frac{27}{16}\mathsf S(6)^6.
\]
For a smooth solenoidal approximation, the \(L^2\) pairing of the projected
equation with \(-\Delta u\) gives
\begin{equation}
 \frac12Y'(t)+\nu Z(t)
 =\operatorname{Re}\ip{(u\cdot\nabla)u}{\Delta u}_{L^2}
 -\operatorname{Re}\ip{g}{\Delta u}_{L^2}.
 \label{f:eq:qc-H1-energy-pairing}
\end{equation}
\begin{samepage}
The explicit three-dimensional Sobolev inequality and interpolation give
\[
 \norm u_6\le\mathsf S(6)Y^{1/2},
 \qquad
 \norm{\nabla u}_3
 \le\mathsf S(6)^{1/2}Y^{1/4}Z^{1/4}.
\]
Consequently,
\[
 \abs{\ip{(u\cdot\nabla)u}{\Delta u}_{L^2}}
 \le \mathsf S(6)^{3/2}Y^{3/4}Z^{3/4}
 \le\frac\nu2Z+\frac{27\mathsf S(6)^6}{32\nu^3}Y^3,
\]
while
\[
 \abs{\ip{g}{\Delta u}_{L^2}}
 \le\frac\nu4Z+\nu^{-1}\norm g_2^2.
\]
Consequently,
\begin{equation}
 Y'(t)+\frac\nu2 Z(t)
 \le c_1\nu^{-3}Y(t)^3+2\nu^{-1}\norm{g(t)}_2^2.
 \label{f:eq:qc-H1-ode}
\end{equation}
\par
\end{samepage}
\begin{samepage}
The integral inequality in \eqref{f:eq:qc-H1-ode}, with the force contribution
kept in the same bootstrap, shows that
\begin{equation}
 Y(t)\le2R_{a,f}
 \label{f:eq:qc-H1-local-explicit}
\end{equation}
whenever \(0\le t\le1\) and
\(8c_1\nu^{-3}R_{a,f}^2t\le1\).
\par
\end{samepage}
\begin{samepage}
The estimate holds on the explicit lifespan
\begin{equation}
 \begin{aligned}
 \tau_{\nu,f}(a):=
 \min\left\{1,
 \frac{\nu^3}{8c_1R_{a,f}^2}
 \right\},\\
 \sup_{0\le t\le\tau_{\nu,f}(a)}Y(t)\le2R_{a,f},
 \qquad
 \int_0^{\tau_{\nu,f}(a)}Z(t)\dd t<\infty.
 \end{aligned}
 \label{f:eq:qc-H1-local-lifespan}
\end{equation}
\par
\end{samepage}
\begin{samepage}
Let \(\mathcal J_k\) be the orthogonal Fourier projection to
\(\{\abs\xi\le k\}\).  On the range of \(\mathcal J_k\), the regularized
problem
\begin{equation}
 \partial_tu^{(k)}-\nu\Delta u^{(k)}
 =-\mathcal J_k\Pp\nabla\!\cdot
   (u^{(k)}\otimes u^{(k)})+\mathcal J_kg,
 \qquad u^{(k)}(0)=\mathcal J_ka,
 \label{f:eq:qc-H1-Friedrichs-system}
\end{equation}
is a Banach-space ordinary differential equation in \(H^1_\sigma\).
\end{samepage}
\begin{samepage}
Indeed, if
\[
 \mathcal N_k(z):=\mathcal J_k\Pp\nabla\!\cdot(z\otimes z),
\]
then the Bernstein inequalities on the range of \(\mathcal J_k\) give
\begin{equation}
 \norm{\mathcal N_k(z)-\mathcal N_k(y)}_{H^1}
 \le C_k\bigl(\norm z_{H^1}+\norm y_{H^1}\bigr)
       \norm{z-y}_{H^1},
 \qquad C_k<\infty.
 \label{f:eq:qc-H1-Friedrichs-Lipschitz}
\end{equation}
\end{samepage}
Picard's theorem supplies a maximal trajectory.  The projection is
self-adjoint and commutes with \(\Pp\), \(\nabla\), and \(\Delta\), so the
forced kinetic identity holds along that trajectory.  Since
\begin{equation}
 \norm z_{H^1}^2\le(1+k^2)\norm z_2^2
 \qquad(z=\mathcal J_kz),
 \label{f:eq:qc-H1-Friedrichs-band-control}
\end{equation}
the forced kinetic bound \(\norm{u^{(k)}(t)}_2\le
\norm a_2+\int_0^t\norm{g(s)}_2\dd s\) and the Banach-space continuation criterion extend the
trajectory for all positive times.  Estimate
\eqref{f:eq:qc-H1-ode} holds for \(u^{(k)}\), with constants independent of
\(k\).  Equations \eqref{f:eq:qc-H1-local-explicit}--%
\eqref{f:eq:qc-H1-local-lifespan} give
\begin{equation}
 \sup_k\left(
 \norm{u^{(k)}}_{L^\infty(0,\tau_{\nu,f}(a);H^1)}
 +\norm{u^{(k)}}_{L^2(0,\tau_{\nu,f}(a);H^2)}
 \right)<\infty.
 \label{f:eq:qc-H1-Friedrichs-uniform}
\end{equation}
\begin{samepage}
The nonlinear estimate
\[
 \norm{\Pp\nabla\!\cdot(u\otimes u)}_2
 \le \norm u_6\norm{\nabla u}_3
 \le \mathsf S(6)^{3/2}Y^{3/4}Z^{1/4},
\]
shows that the regularized nonlinear terms are bounded in
\(L^2(0,\tau_{\nu,f}(a);L^2)\), uniformly in \(k\).
\par
\end{samepage}
The equations and \(\mathcal J_kg\in L^2(0,\tau_{\nu,f}(a);L^2)\) bound
\(\partial_tu^{(k)}\) in this space.
\begin{samepage}
Weak compactness and the
Aubin--Lions--Simon lemma \cite{Aubin1963,Simon1987} on every ball give, after diagonal
extraction,
\begin{align*}
 u^{(k)}&\stackrel{*}{\rightharpoonup} u
 &&\text{in }L^\infty(0,\tau_{\nu,f}(a);H^1),\\
 u^{(k)}&\rightharpoonup u
 &&\text{in }L^2(0,\tau_{\nu,f}(a);H^2),\\
 u^{(k)}&\longrightarrow u
 &&\text{in }L^2(0,\tau_{\nu,f}(a);H^1(B_R))
 \quad(R<\infty).
\end{align*}
\par
\end{samepage}
The last convergence passes the quadratic term in distributions; lower
semicontinuity preserves \eqref{f:eq:qc-H1-Friedrichs-uniform}.  The displayed
nonlinear estimate now gives
\(\Pp\nabla\!\cdot(u\otimes u)\in L^2L^2\), and the equation gives
\(u_t\in L^2(0,\tau_{\nu,f}(a);L^2)\).
\begin{samepage}
For every
\(\phi\in C_c^\infty(\R^3;\R^3)\) with \(\nabla\cdot\phi=0\) and every
\(\eta\in C_c^\infty([0,\tau_{\nu,f}(a)))\), passage to the limit in the
time-integrated regularized equation gives
\begin{align}
 -\int_0^{\tau_{\nu,f}(a)}\ip{u}{\phi}_{L^2}\eta'\dd t
 &+\nu\int_0^{\tau_{\nu,f}(a)}\ip{\nabla u}{\nabla\phi}_{L^2}\eta\dd t
 -\int_0^{\tau_{\nu,f}(a)}\ip{u\otimes u}{\nabla\phi}_{L^2}\eta\dd t
 \notag\\
 &=\ip{a}{\phi}_{L^2}\eta(0)
 +\int_0^{\tau_{\nu,f}(a)}\ip{g}{\phi}_{L^2}\eta\dd t.
 \label{f:eq:qc-H1-Friedrichs-initial-trace}
\end{align}
The right-hand side fixes the weak initial trace as \(a\).
\par
\end{samepage}
The Lions--Magenes theorem \cite{LionsMagenes1972} for
\(H^2_\sigma\hookrightarrow H^1_\sigma\hookrightarrow L^2_\sigma\)
gives \(u\in C([0,\tau_{\nu,f}(a)];H^1_\sigma)\); its continuous trace at zero is
the trace in \eqref{f:eq:qc-H1-Friedrichs-initial-trace}, hence \(u(0)=a\).
Variation of constants gives the mild equation, and
\(p=R_iR_j(u_i u_j)-(-\Delta)^{-1}\nabla\cdot f\) supplies its normalized
pressure.  Since
\(u_t\in L^2L^2\) and \(u\in L^2H^2\), pairing the limiting equation with
\(u\) is legitimate; incompressibility cancels the transport term and proves
the forced kinetic-energy identity.

If \(u,v\) are two such solutions and \(w=u-v\), incompressibility cancels
the \((u\cdot\nabla)w\) pairing and gives
\[
 \frac12\frac{\dd}{\dd t}\norm w_2^2
 +\nu\norm{\nabla w}_2^2
 =-\operatorname{Re}\ip{(w\cdot\nabla)v}{w}_{L^2}.
\]
This Sobolev constant yields
\[
 \abs{\ip{(w\cdot\nabla)v}{w}_{L^2}}
 \le \mathsf S(6)^{3/2}\norm{\nabla v}_2
       \norm w_2^{1/2}\norm{\nabla w}_2^{3/2}
 \le\frac\nu2\norm{\nabla w}_2^2
 +\frac{27\mathsf S(6)^6}{32\nu^3}
  \norm{\nabla v}_2^4\norm w_2^2.
\]
Gronwall's inequality proves uniqueness.  This calculation one derivative
higher also gives the finite-interval stability needed below.  Indeed, pairing
the difference equation with \(-\Delta w\), integrating the
\((v\cdot\nabla)w\) term once in space, and applying Young's inequality gives
\begin{equation}
 \frac{\dd}{\dd t}\norm{\nabla w}_2^2
 +\nu\norm{\Delta w}_2^2
 \le \frac{C_{\rm st}}\nu
 \left(\norm{\nabla u}_3^2+\norm{\nabla v}_3^2\right)
 \norm{\nabla w}_2^2.
 \label{f:eq:qc-H1-stability}
\end{equation}
The coefficient is integrable because
\[
 \norm{\nabla u}_3^2
 \le \mathsf S(6)\norm{\nabla u}_2\norm{\Delta u}_2,
\]
and likewise for \(v\).  The \(L^2\) difference estimate and
\eqref{f:eq:qc-H1-stability} show continuous dependence in
\(C([0,S];H^1)\) on every finite interval \([0,S]\).

Let \(c_0>0\) be the absolute smallness constant in the full-space
\(H^1\) local existence theorem, Theorem~5.4(ii) of \cite{Tao2013}.
Applied on an interval \(I\) of length \(h\), with left trace \(b\) and
projected source \(g\), its viscosity-one hypothesis is
\[
 h\left(\norm b_{H^1}+\norm g_{L^1(I;H^1)}\right)^4\le c_0.
\]
After viscosity rescaling, the hypothesis becomes
\begin{equation}
 h\left(\norm b_{H^1}+\norm g_{L^1(I;H^1)}\right)^4
 \le c_0\nu^3.
 \label{f:eq:qc-H1-smoothing-smallness}
\end{equation}
Fix \(0<\varepsilon<S\le\tau_{\nu,f}(a)\), and put
\[
 K_{S,f}:=1+\sup_{0\le t\le S}\norm{u(t)}_{H^1}
 +\sup_{0\le t\le S}\norm{f(t)}_{H^2}.
\]
Choose
\[
 0<h\le\min\left\{\varepsilon,
 c_0\nu^3K_{S,f}^{-4}\right\}.
\]
For \(I=[0,h]\) or \(I=[t-h,t]\), \(h\le t\le S\), and left endpoint
\(s\), the facts that \(S\le\tau_{\nu,f}(a)\le1\) and
\(g=\Pp f\) give
\[
 \norm{u(s)}_{H^1}+\norm g_{L^1(I;H^1)}
 \le \sup_{0\le r\le S}\norm{u(r)}_{H^1}
 +h\sup_{0\le r\le S}\norm{f(r)}_{H^2}
 \le K_{S,f}.
\]
Thus these intervals obey \eqref{f:eq:qc-H1-smoothing-smallness}.  The
quantitative positive-time regularity estimate in Lemma~5.5 of
\cite{Tao2013}, after viscosity rescaling with interior time \(h/2\), applies
to this inhomogeneous mild formula.  Its source terms are bounded by the
displayed smooth-force norms, so the heat-semigroup/product iteration gives,
for every \(m\in\N\), with
\[
 \Gamma_{m,S}(f):=
 \sum_{j=0}^{m}\sup_{0\le t\le S}
 \norm{\partial_t^jf(t)}_{H^{m+2}},
\]
\begin{equation}
 \sup_{\varepsilon\le t\le S}\norm{u(t)}_{H^m}
 \le C_{m,\varepsilon,S}
 \bigl(\nu,K_{S,f},\Gamma_{m,S}(f)\bigr)<\infty.
 \label{f:eq:qc-H1-positive-time-smoothing}
\end{equation}
Apply the mild formula from time \(\varepsilon/2\).  The bounds
\eqref{f:eq:qc-H1-positive-time-smoothing}, the heat multiplier, and
Lemma~\ref{f:lem:sobolev-product} show that
\(u\in C([\varepsilon,S];H^m)\) for every \(m\).  The identities
\begin{equation}
 p=R_iR_j(u_i u_j)-(-\Delta)^{-1}\nabla\cdot f,
 \qquad
 u_t=\nu\Delta u-\Pp\nabla\!\cdot(u\otimes u)+g
 \label{f:eq:qc-H1-spacetime-bootstrap}
\end{equation}
then give
\(p\in C([\varepsilon,S];H^m)\) and
\(u_t\in C([\varepsilon,S];H^m)\) for every \(m\).  Differentiating
\eqref{f:eq:qc-H1-spacetime-bootstrap} and the pressure identity, and applying
this product estimate, proves inductively that all temporal derivatives
of \(u\) and \(p\) belong to every spatial Sobolev space continuously on
\([\varepsilon,S]\).  Sobolev embedding proves
\((u,p)\in C^\infty((0,\tau_{\nu,f}(a)]\times\R^3)\).

Choose \(0<t_0<\tau_{\nu,f}(a)\).  Equation
\eqref{f:eq:qc-H1-positive-time-smoothing} places \(u(t_0)\) in
\(H^\infty_\sigma\).  The absolute-time form of
Theorem~\ref{f:thm:ier-global-smoothness}, begun from \(u(t_0)\) at \(t_0\),
supplies a global smooth history driven for \(t\ge t_0\) by the original
prescribed force \(f(t)\).  The
uniqueness estimate above identifies it with the local history on their
overlap, so the two histories glue to \eqref{f:eq:qc-global-H1-class}.
This estimate proves uniqueness in the full displayed class.  The forced
kinetic identity on the first interval and the classical forced identity
after \(t_0\) agree
at the common trace and prove the final assertion.
\end{proof}

\subsection{Finite norm and trace majorants}

The compatible rectangles carry more information than the endpoint argument
uses.  We first name scalar majorants of their finite norms and close the estimates that will
feed the mixed rows, frequency tails, and continuation bounds below.

\begin{definition}[Rectangle-norm and trace majorants]
\label{f:def:qc-producer-constants}
Let \(u_0\in\HsSigma(\R^3)\), \(\nu>0\), prescribe
\(f\in C^\infty([0,\infty);\mathcal S(\R^3;\R^3))\), and let \((u,p)\) be
the pressure-normalized solution generated by \((u_0,\nu,f)\).  Let
\(0<T<\infty\) satisfy
either \(T<T_*\) or \(T=T_*<\infty\).  Set
\(I_T=(0,T)\),
\[
 V_n:=\partial_t^n u,
 \qquad
 Q_n:=\partial_t^n p,
 \qquad
 F_n:=\partial_t^n f,
 \qquad
 \Phi_n:=(-\Delta)^{-1}\nabla\cdot F_n,
 \qquad n\in\Nzero.
\]
Define the monotone rectangle-norm majorant
\begin{equation}
 \mathfrak P_{N,M}[u_0,\nu,f;T]
 :=1+\sup_{0<S<T}\mathfrak R_{N,M}(u;S),
 \label{f:eq:qc-producer-constant}
\end{equation}
and the trace constant
\begin{equation}
 \mathfrak T_{N,M}[u_0,\nu,f;T]
 :=\left(1+T^{-1}\right)^{1/2}
   \mathfrak P_{N,M}[u_0,\nu,f;T]^{1/2}.
 \label{f:eq:qc-trace-constant}
\end{equation}
All occurrences of \(\mathfrak P_{N,M}\) and \(\mathfrak T_{N,M}\) below
carry the arguments \([u_0,\nu,f;T]\).
The velocity datum \(u_0\), viscosity \(\nu\), prescribed force \(f\),
pressure-normalized solution \((u,p)\), horizon \(T\), interval \(I_T\),
and row notation \(V_n,Q_n,F_n,\Phi_n\) fixed here remain in force
throughout the rest of this section.
In the terminal alternative \(T=T_*<\infty\), the full-interval representatives
are those of
Proposition~\ref{f:prop:datum-terminal-realization}, and every endpoint value
used below is the compatible trace supplied by
Theorem~\ref{f:thm:ier-compatible-endpoint-traces}.
\end{definition}

\begin{lemma}[Majorant finiteness]
\label{f:lem:qc-producer-finiteness}
For every \(N,M\in\Nzero\),
\begin{equation}
 1\le \mathfrak P_{N,M}[u_0,\nu,f;T]<\infty.
 \label{f:eq:qc-producer-finiteness}
\end{equation}
\end{lemma}

\begin{proof}
If \(T=T_*<\infty\), equation~\eqref{f:eq:datum-rectangle-bound} proves
\eqref{f:eq:qc-producer-finiteness}.  If \(T<T_*\), classical persistence on
\([0,T]\) gives
\[
 \sup_{0\le t\le T}\norm{V_n(t)}_{H^m}<\infty
 \qquad(n,m\in\Nzero),
\]
and \eqref{f:eq:finite-rectangle-norm} proves
\eqref{f:eq:qc-producer-finiteness}.
\end{proof}

\Needspace{12\baselineskip}
\begin{lemma}[Fixed analytic constants]
\label{f:lem:qc-analytic-constants}
For \(r\in\Nzero\), set
\[
 A_r:=2C_r^{\rm alg}.
\]
Then
\begin{equation}
 \norm{R_iR_j(f_i g_j)}_{H^r}
 +\norm{\Pp\nabla\!\cdot(f\otimes g)}_{H^{r-1}}
 \le A_r\norm f_{H^{r+1}}\norm g_{H^{r+1}}.
 \label{f:eq:qc-product-constant}
\end{equation}
This constant satisfies
\begin{equation}
 \norm{R_iR_j(f_i g_j)}_{H^r}
 +\norm{\nabla R_iR_j(f_i g_j)}_{H^{r-1}}
 \le A_r\norm f_{H^{r+1}}\norm g_{H^{r+1}}.
 \label{f:eq:qc-pressure-product-constant}
\end{equation}
\begin{samepage}
For \(s,k\in\Nzero\), \(2\le b\le\infty\), and
\(s>k+\frac32-\frac3b\), define
\[
 \nabla^k f:=(\partial^\alpha f)_{|\alpha|=k},
 \qquad
 \norm{\nabla^k f}_{L^b}
 :=\left\|\left(
   \sum_{|\alpha|=k}|\partial^\alpha f|^2
 \right)^{1/2}\right\|_{L^b},
\]
and set
\[
 d_k:=\#\{\alpha\in\Nzero^3:\abs\alpha=k\},
 \qquad
 \frac1{q_b}:=\frac12-\frac1b,
\]
with \(q_2=\infty\), and
\begin{equation}
 S_{s,k,b}:=
 d_k^{1/2}(2\pi)^{-3(1/2-1/b)}
 \norm{\abs\xi^k\langle\xi\rangle^{-s}}_{L^{q_b}_\xi}.
 \label{f:eq:qc-embedding-constant-definition}
\end{equation}
Then \(S_{s,k,b}<\infty\) and
\begin{equation}
 \norm{\nabla^k f}_{L^b}\le S_{s,k,b}\norm f_{H^s}.
 \label{f:eq:qc-embedding-constant}
\end{equation}
\par
\end{samepage}
For \(m\ge2\), set \(L_m:=\max\{1,C_m^{\rm alg}\}\).  Then
\begin{align}
 \norm{fg}_{H^m}
 &\le L_m\norm f_{H^m}\norm g_{H^m},
 \label{f:eq:qc-algebra-constant}\\
 \norm{e^{\nu t\Delta}\Pp\nabla\!\cdot F}_{H^m}
 &\le L_m(\nu t)^{-1/2}\norm F_{H^m},
 \qquad t>0.
 \label{f:eq:qc-heat-constant}
\end{align}
\end{lemma}

\begin{proof}
Lemmas~\ref{f:lem:sobolev-product} and \ref{f:lem:riesz-leray}, applied
componentwise before summing squares, give
\begin{align*}
 \norm{R_iR_j(f_i g_j)}_{H^r}
 &\le\norm{f\otimes g}_{H^r}
 \le C_r^{\rm alg}\norm f_{H^{r+1}}\norm g_{H^{r+1}},\\
 \norm{\Pp\nabla\!\cdot(f\otimes g)}_{H^{r-1}}
 &\le\norm{f\otimes g}_{H^r}
 \le C_r^{\rm alg}\norm f_{H^{r+1}}\norm g_{H^{r+1}}.
\end{align*}
Moreover,
\[
 \norm{\nabla R_iR_j(f_i g_j)}_{H^{r-1}}
 \le\norm{R_iR_j(f_i g_j)}_{H^r}.
\]
This proves \eqref{f:eq:qc-product-constant} and
\eqref{f:eq:qc-pressure-product-constant}.

Let \(b'=b/(b-1)\).  Interpolation between Fourier inversion
\(L^1_\xi\to L^\infty_x\), with norm \((2\pi)^{-3/2}\), and Plancherel
gives
\[
 \norm{\mathcal F^{-1}h}_{L^b_x}
 \le(2\pi)^{-3(1/2-1/b)}\norm h_{L^{b'}_\xi}.
\]
For \(\abs\alpha=k\), H\"older's inequality with
\(1/b'=1/2+1/q_b\) gives
\[
 \norm{\partial^\alpha f}_{L^b}
 \le(2\pi)^{-3(1/2-1/b)}
 \norm{\abs\xi^k\langle\xi\rangle^{-s}}_{L^{q_b}}
 \norm f_{H^s}.
\]
The weight belongs to \(L^{q_b}\) under the displayed strict inequality
for \(s\); for \(b=2\), its \(L^\infty\)-norm is finite because \(s>k\).
Summation over \(\abs\alpha=k\) proves
\eqref{f:eq:qc-embedding-constant}.

The first estimate in \eqref{f:eq:qc-algebra-constant} is
Lemma~\ref{f:lem:sobolev-product}.  For the heat estimate, the symbol of
\(\Pp\nabla\!\cdot\), acting from tensors to vectors, has norm at most
\(\abs\xi\), while
\[
 \sup_{\rho\ge0}\rho e^{-\nu t\rho^2}=(2e\nu t)^{-1/2}.
\]
Plancherel gives
\[
 \norm{e^{\nu t\Delta}\Pp\nabla\!\cdot F}_{H^m}
 \le(2e\nu t)^{-1/2}\norm F_{H^m}
 \le L_m(\nu t)^{-1/2}\norm F_{H^m}.
\]
\end{proof}

\begin{definition}[Canonical rectangle majorants]
\label{f:def:qc-canonical-majorants}
For \(N,M\in\Nzero\), define
\begin{equation}
 \mathfrak B_{N,M}[u_0,\nu,f;T]
 :=\inf\left\{B\ge0:
 \mathfrak R_{N,M}(u;S)\le B
 \text{ for every }0<S<T\right\}.
 \label{f:eq:qc-canonical-rectangle-majorant}
\end{equation}
Monotone convergence gives
\begin{equation}
 \mathfrak B_{N,M}[u_0,\nu,f;T]
 =\begin{cases}
   \mathfrak R_{N,M}(u;T),&T<T_*,\\
   \mathfrak R_{N,M}^*[u_0,\nu,f,T_*],&T=T_*<\infty.
  \end{cases}
 \label{f:eq:qc-canonical-rectangle-identification}
\end{equation}
The pressure-normalized history in
Definition~\ref{f:def:qc-producer-constants} is fixed for \((u_0,\nu,f)\).
Thus
\(\mathfrak B_{N,M}[u_0,\nu,f;T]\) is the exact minimal rectangle bound of
that datum-generated history, not a supremum over a family of solutions and
not an independently chosen constant.
Equations~\eqref{f:eq:rectangle-matrix-norm-identities} and
\eqref{f:eq:qc-canonical-rectangle-majorant} give the exact matrix identities
\begin{align}
 \mathfrak B_{N,M}[u_0,\nu,f;T]
 &=\sup_{0<S<T}\norm{\mathbf G_{N,M}(S)}_{\mathrm F}^2,
 \label{f:eq:qc-B-matrix-identity}\\
 \mathfrak P_{N,M}[u_0,\nu,f;T]
 &=1+\sup_{0<S<T}\norm{\mathbf G_{N,M}(S)}_{\mathrm F}^2.
 \label{f:eq:qc-P-matrix-identity}
\end{align}
Consequently,
\begin{equation}
 \mathfrak P_{N,M}[u_0,\nu,f;T]
 =1+\mathfrak B_{N,M}[u_0,\nu,f;T].
 \label{f:eq:qc-producer-through-B}
\end{equation}

For \(n\in\Nzero\) and \(m\in\mathbb Z\), define
\begin{align}
 \mathfrak G^{(0)}_{n,m}
 &:=\norm{F_n(0)}_{H^m},
 &
 \mathfrak G^{(1)}_{n,m}
 &:=\norm{F_n}_{L^1(I_T;H^m)},
 \notag\\
 \mathfrak G^{(2)}_{n,m}
 &:=\norm{F_n}_{L^2(I_T;H^m)},
 &
 \mathfrak G^{(\infty)}_{n,m}
 &:=\sup_{0\le t\le T}\norm{F_n(t)}_{H^m},
 \label{f:eq:qc-force-row-majorants}
\end{align}
For \(n,m\in\Nzero\), define the force-pressure coordinates
\begin{align}
 \mathfrak L^{(1)}_{n,m}
 &:=\norm{\Phi_n}_{L^1(I_T;H^m)}
   +\norm{\nabla\Phi_n}_{L^1(I_T;H^{m-1})},
 &
 \mathfrak L^{(2)}_{n,m}
 &:=\norm{\Phi_n}_{L^2(I_T;H^m)},
 \notag\\
 \mathfrak L^{(\infty)}_{n,m}
 &:=\sup_{0\le t\le T}\left(
 \norm{\Phi_n(t)}_{H^m}
 +\norm{\nabla\Phi_n(t)}_{H^{m-1}}\right).
 \label{f:eq:qc-force-pressure-majorants}
\end{align}
Every one of these coordinates is finite by the prescribed force class and
\eqref{f:eq:force-pressure-potential-bound}.  They belong to the same forced
history and are not independently chosen sources.

For \(n,m\in\Nzero\), define the initial-jet majorants
\(\mathfrak D_{n,m}=\mathfrak D_{n,m}[u_0,\nu,f]\) by
\begin{align}
 \mathfrak D_{0,m}&:=\norm{u_0}_{H^m},
 \label{f:eq:qc-D0}\\
 \mathfrak D_{n+1,m}
 &:=\nu\mathfrak D_{n,m+2}
 +A_{m+1}\sum_{a=0}^{n}\binom na
   \mathfrak D_{a,m+2}\mathfrak D_{n-a,m+2}
  +\mathfrak G^{(0)}_{n,m}.
 \label{f:eq:qc-Drec}
\end{align}
\begin{samepage}
Define the pointwise recursion
\(\mathfrak Y_{n,m}=\mathfrak Y_{n,m}[u_0,\nu,f;T]\) by
\begin{align}
 \mathfrak Y_{0,0}&:=K_f(T),
 \label{f:eq:qc-Y00}\\
 \mathfrak Y_{0,m}
 &:=\min\left\{
 (1+T^{-1})^{1/2}\mathfrak B_{0,m}^{1/2},
 \norm{u_0}_{H^m}+T^{1/2}\mathfrak B_{0,m+1}^{1/2}
 \right\},
 \qquad m\ge1,
 \label{f:eq:qc-Y0}\\
 \mathfrak Y_{n+1,m}
 &:=\nu\mathfrak Y_{n,m+2}
 +A_{m+1}\sum_{a=0}^{n}\binom na
 \mathfrak Y_{a,m+2}\mathfrak Y_{n-a,m+2}
 +\mathfrak G^{(\infty)}_{n,m}.
 \label{f:eq:qc-Yrec}
\end{align}
\end{samepage}
\begin{samepage}
For \(m\in\{-1\}\cup\Nzero\), put
\[
 m_+:=\max\{m,0\},
 \qquad
 \mu(m):=\max\{m-1,0\},
\]
and define
\begin{equation}
 \widetilde{\mathfrak X}_{n,m}
 :=\min\left\{
 T^{1/2}\mathfrak Y_{n,m_+},
 \mathfrak B_{n,\mu(m)}^{1/2}
 \right\}.
 \label{f:eq:qc-Xraw}
\end{equation}
\end{samepage}
Set
\begin{align}
 \mathfrak X_{0,0}
 &:=\min\left\{\widetilde{\mathfrak X}_{0,0},
 T^{1/2}K_f(T)\right\},
 \label{f:eq:qc-X00}\\
 \mathfrak X_{0,1}
 &:=\min\left\{\widetilde{\mathfrak X}_{0,1},
 \left(T+\frac1{2\nu}\right)^{1/2}K_f(T)\right\},
 \label{f:eq:qc-X01}\\
 \mathfrak X_{n,m}&:=\widetilde{\mathfrak X}_{n,m}
 \quad\text{for all other }(n,m).
 \label{f:eq:qc-Xotherwise}
\end{align}
For \(n,m\in\Nzero\), define
\begin{equation}
 \mathfrak Z_{n,m}[u_0,\nu,f;T]
 :=\min\left\{
 \mathfrak Y_{n,m},
 (1+T^{-1})^{1/2}\mathfrak B_{n,m}^{1/2},
 \mathfrak D_{n,m}+T^{1/2}\mathfrak X_{n+1,m}
 \right\}.
 \label{f:eq:qc-Z}
\end{equation}
All \(\mathfrak B,\mathfrak D,\mathfrak X,\mathfrak Y,\mathfrak Z\),
\(\mathfrak G\), and \(\mathfrak L\) in
this section carry the arguments \([u_0,\nu,f;T]\), except that
\(\mathfrak D\) is independent of \(T\).
\end{definition}

\begin{proposition}[Exact viscosity and horizon scaling]
\label{f:prop:qc-exact-viscosity-scaling}
Let
\[
 a:=\nu^{-1}u_0,
 \qquad
 v(s,x):=\nu^{-1}u(s/\nu,x),
 \qquad
 q(s,x):=\nu^{-2}p(s/\nu,x),
 \qquad
 \widetilde f(s,x):=\nu^{-2}f(s/\nu,x).
\]
Then \((v,q)\) is the unit-viscosity solution generated by
\((a,1,\widetilde f)\).  If
\[
 \widetilde V_n:=\partial_s^nv,
 \qquad
 \widetilde Q_n:=\partial_s^nq,
 \qquad
 \widetilde F_n:=\partial_s^n\widetilde f,
\]
then
\begin{equation}
 V_n(t)=\nu^{n+1}\widetilde V_n(\nu t),
 \qquad
 Q_n(t)=\nu^{n+2}\widetilde Q_n(\nu t),
 \qquad
 F_n(t)=\nu^{n+2}\widetilde F_n(\nu t).
 \label{f:eq:qc-row-viscosity-scaling}
\end{equation}
Substitution of \eqref{f:eq:qc-row-viscosity-scaling} into the defining norms gives
\begin{align}
 \mathfrak B_{N,M}[u_0,\nu,f;T]
 =\sum_{n=0}^{N}\Bigl[{}
 &\nu^{2n+1}
 \norm{\widetilde V_n}_{L^2(0,\nu T;H^{M+1})}^2
 \notag\\[-2pt]
 &+\nu^{2n+3}
 \norm{\widetilde V_{n+1}}_{L^2(0,\nu T;H^{M-1})}^2
 \Bigr],
 \label{f:eq:qc-rectangle-viscosity-scaling}\\
 \norm{Q_n}_{L^1(0,T;H^m)}
 +\norm{\nabla Q_n}_{L^1(0,T;H^{m-1})}
 ={}&\nu^{n+1}\Bigl(
 \norm{\widetilde Q_n}_{L^1(0,\nu T;H^m)}
 \notag\\[-2pt]
 &\hspace{32pt}+
 \norm{\nabla\widetilde Q_n}_{L^1(0,\nu T;H^{m-1})}
 \Bigr),
 \label{f:eq:qc-pressure-viscosity-scaling}\\
 \sup_{0\le t\le T}\norm{Q_n(t)}_{H^m}
 ={}&\nu^{n+2}
 \sup_{0\le s\le\nu T}\norm{\widetilde Q_n(s)}_{H^m}.
 \label{f:eq:qc-pressure-sup-viscosity-scaling}
\end{align}
The initial-jet recursion obeys
\begin{equation}
 \mathfrak D_{n,m}[u_0,\nu,f]
 =\nu^{n+1}\mathfrak D_{n,m}[u_0/\nu,1,\widetilde f].
 \label{f:eq:qc-initial-jet-viscosity-scaling}
\end{equation}
Write
\[
 E_\sigma^u(t):=\frac12\norm{\Lambda^\sigma u(t)}_2^2,
 \qquad
 E_\sigma^v(s):=\frac12\norm{\Lambda^\sigma v(s)}_2^2.
\]
For every \(j\in\Nzero\), the critical-height rows obey
\begin{align}
 E_{j+1/2}^u(t)
 &=\nu^2E_{j+1/2}^v(\nu t),
 \label{f:eq:qc-critical-height-viscosity-scaling}\\
 \int_0^T E_{j+1/2}^u(t)\dd t
 &=\nu\int_0^{\nu T}E_{j+1/2}^v(s)\dd s.
 \label{f:eq:qc-critical-height-mass-viscosity-scaling}
\end{align}
For every \(t_0\ge0\) and \(b\in H^3_\sigma(\R^3)\), put
\[
 f_{\nu,t_0}(\sigma,x):=\nu^{-2}f(t_0+\sigma/\nu,x).
\]
The maximal elapsed lifespan in the shifted mild class has the exact
covariance
\begin{equation}
 \tau_{\nu,f,3}^{\max}(t_0,b)
 =\nu^{-1}\tau_{1,f_{\nu,t_0},3}^{\max}(0,b/\nu).
 \label{f:eq:qc-maximal-lifespan-viscosity-scaling}
\end{equation}
\end{proposition}

\begin{proof}
The change of variables gives
\(u(t)=\nu v(\nu t)\) and \(p(t)=\nu^2q(\nu t)\).
Differentiation proves \eqref{f:eq:qc-row-viscosity-scaling}; the substitution
\(s=\nu t\) proves
\eqref{f:eq:qc-rectangle-viscosity-scaling}--%
\eqref{f:eq:qc-pressure-sup-viscosity-scaling}.
The force-row identity in \eqref{f:eq:qc-row-viscosity-scaling} makes the
force term in \eqref{f:eq:qc-Drec} homogeneous of the same degree as
\(V_{n+1}\).  Equation \eqref{f:eq:qc-initial-jet-viscosity-scaling} then
follows by induction in \eqref{f:eq:qc-D0}--\eqref{f:eq:qc-Drec}.  The amplitude relation
\(u(t)=\nu v(\nu t)\) proves
\eqref{f:eq:qc-critical-height-viscosity-scaling}, and the substitution
\(s=\nu t\) proves
\eqref{f:eq:qc-critical-height-mass-viscosity-scaling}.  For
\((w,\pi)\in\mathscr M_{\nu,f}^3(t_0,t_0+\delta;b)\), set
\[
 \mathcal S_\nu(w,\pi)(\sigma)
 :=
 \left(
  \nu^{-1}w(t_0+\sigma/\nu),
  \nu^{-2}\pi(t_0+\sigma/\nu)
 \right),
 \qquad 0\le\sigma\le\nu\delta.
\]
Time translation and this change of variables show that
\(\mathcal S_\nu\) is a bijection from
\(\mathscr M_{\nu,f}^3(t_0,t_0+\delta;b)\) onto
\(\mathscr M_{1,f_{\nu,t_0}}^3(0,\nu\delta;b/\nu)\).
Its inverse is the corresponding amplitude and time rescaling,
which proves \eqref{f:eq:qc-maximal-lifespan-viscosity-scaling}.
\end{proof}

\begin{proposition}[Height-coordinate reconstruction of the finite rows]
\label{f:prop:qc-critical-height-row-propagation}
For \(m\in\Nzero\), define
\begin{align}
 \mathfrak H_m[u_0,\nu,f;T]
 &:={}
 \frac12\norm{\Lambda^{m+1/2}u}_{L^2(I_T;L^2)}^2,
 \label{f:eq:qc-height-producer-coordinate}\\
 \mathfrak S_m[u_0,\nu,f;T]^2
 &:={}
 2^m\left(TK_f(T)^2+2\mathfrak H_m[u_0,\nu,f;T]\right).
 \label{f:eq:qc-spatial-row-from-height}
\end{align}
Then
\begin{equation}
 2\mathfrak H_m[u_0,\nu,f;T]
 \le\mathfrak B_{0,m}[u_0,\nu,f;T],
 \qquad
 \norm u_{L^2(I_T;H^m)}\le\mathfrak S_m[u_0,\nu,f;T].
 \label{f:eq:qc-height-spatial-bounds}
\end{equation}
Define recursively, for \(n,m\in\Nzero\),
\begin{align}
 \mathfrak A^{\rm v}_{0,m}
 &:={}
 T^{-1/2}\mathfrak S_m
 +\nu T^{1/2}\mathfrak S_{m+2}
 +A_{m+1}\mathfrak S_{m+2}^2
 +\mathfrak G^{(1)}_{0,m},
 \label{f:eq:qc-height-base-pointwise-majorant}\\
 \mathfrak A^{\rm p}_{n,m}
 &:={}
 A_m\sum_{a=0}^{n}\binom na
 \mathfrak A^{\rm v}_{a,m+1}
 \mathfrak A^{\rm v}_{n-a,m+1}
 +\mathfrak L^{(\infty)}_{n,m},
 \label{f:eq:qc-height-pressure-majorant}\\
 \mathfrak A^{\rm v}_{n+1,m}
 &:={}
 \nu\mathfrak A^{\rm v}_{n,m+2}
 +A_{m+1}\sum_{a=0}^{n}\binom na
 \mathfrak A^{\rm v}_{a,m+2}
 \mathfrak A^{\rm v}_{n-a,m+2}
 +\mathfrak G^{(\infty)}_{n,m}.
 \label{f:eq:qc-height-temporal-majorant}
\end{align}
Every finite quantity in this recursion is finite, and
\begin{align}
 \sup_{0\le t\le T}\norm{V_n(t)}_{H^m}
 &\le\mathfrak A^{\rm v}_{n,m},
 \label{f:eq:qc-height-temporal-row-bound}\\
 \sup_{0\le t\le T}
 \left(\norm{Q_n(t)}_{H^m}
 +\norm{\nabla Q_n(t)}_{H^{m-1}}\right)
 &\le\mathfrak A^{\rm p}_{n,m}.
 \label{f:eq:qc-height-pressure-row-bound}
\end{align}
Let \(M_-:=\max\{M-1,0\}\) and set
\begin{align}
 \widehat{\mathfrak B}_{N,M}^{\rm h}[u_0,\nu,f;T]
 &:={}
 T\sum_{n=0}^{N}\left(
 (\mathfrak A^{\rm v}_{n,M+1})^2
 +(\mathfrak A^{\rm v}_{n+1,M_-})^2
 \right),
 \label{f:eq:qc-height-reconstructed-rectangle}\\
 \widehat{\mathfrak Q}_{N,M}^{\rm h}[u_0,\nu,f;T]
 &:={}
 T\sum_{n=0}^{N}\mathfrak A^{\rm p}_{n,M}.
 \label{f:eq:qc-height-reconstructed-pressure}
\end{align}
Then
\begin{align}
 \mathfrak B_{N,M}[u_0,\nu,f;T]
 &\le\widehat{\mathfrak B}_{N,M}^{\rm h}[u_0,\nu,f;T],
 \label{f:eq:qc-B-by-height-reconstruction}\\
 \mathfrak Q_{N,M}(p;T)
 &\le\widehat{\mathfrak Q}_{N,M}^{\rm h}[u_0,\nu,f;T].
 \label{f:eq:qc-pressure-by-height-reconstruction}
\end{align}
For fixed \(N,M\), the right-hand sides use only
\begin{equation}
 \begin{gathered}
 \nu,\quad T,\quad K_f(T),\quad N,\quad M,\\
 \{\mathfrak H_q[u_0,\nu,f;T]:0\le q\le M+2N+4\},
 \quad\{A_q:0\le q\le M+2N+4\},\\
 \{\mathfrak G^{(1)}_{0,b}:0\le b\le M+2N+2\},\\
 \{\mathfrak G^{(\infty)}_{a,b}:0\le a\le N,\ 0\le b\le M+2N+2\},
 \quad
 \{\mathfrak L^{(\infty)}_{a,M}:0\le a\le N\}.
 \end{gathered}
 \label{f:eq:qc-height-reconstruction-dependence}
\end{equation}
The quantities \(\mathfrak H_q\) are the exact completed-height coordinates
of the fixed datum-generated history.  Equation
\eqref{f:eq:qc-height-spatial-bounds} converts every finite height into a spatial
Sobolev row, while \eqref{f:eq:qc-B-by-height-reconstruction} uses the forced
equation to recover the temporal and pressure rows of any prescribed finite
rectangle.  Over all finite indices this is the critical-height exhaustion of
the same mixed intersection.
\end{proposition}

\begin{proof}
The first estimate in \eqref{f:eq:qc-height-spatial-bounds} follows from the
upper row of \(\mathfrak B_{0,m}\).  The Fourier inequality
\[
 (1+\abs\xi^2)^m\le2^m(1+\abs\xi^{2m+1})
\]
and \eqref{f:eq:energy-Linf} give the second estimate.

The projected equation and \eqref{f:eq:qc-product-constant} give
\[
 \norm{u_t}_{L^1(I_T;H^m)}
 \le\nu T^{1/2}\mathfrak S_{m+2}
 +A_{m+1}\mathfrak S_{m+2}^2
 +\mathfrak G^{(1)}_{0,m}.
\]
Choose a Lebesgue time \(t_m\in I_T\) satisfying
\[
 \norm{u(t_m)}_{H^m}
 \le T^{-1/2}\norm u_{L^2(I_T;H^m)}.
\]
The Banach-valued fundamental identity gives
\eqref{f:eq:qc-height-temporal-row-bound} at \(n=0\).
Assume that bound through row \(n\).  The forced pressure recurrence and
\eqref{f:eq:qc-pressure-product-constant} give
\eqref{f:eq:qc-height-pressure-row-bound}.  The projected temporal recurrence
with \(\Pp F_n\) retained and \eqref{f:eq:qc-product-constant} give the next velocity row through
\eqref{f:eq:qc-height-temporal-majorant}.  Induction proves
\eqref{f:eq:qc-height-temporal-row-bound}--%
\eqref{f:eq:qc-height-pressure-row-bound}.

Integrating the pointwise bounds over \(I_T\) gives
\eqref{f:eq:qc-B-by-height-reconstruction}; when \(M=0\), use
\(H^0\hookrightarrow H^{-1}\) on the lower row.  This integration in
\eqref{f:eq:qc-height-pressure-row-bound} gives
\eqref{f:eq:qc-pressure-by-height-reconstruction}.  Induction in
\eqref{f:eq:qc-height-temporal-majorant} gives the finite dependency range in
\eqref{f:eq:qc-height-reconstruction-dependence}.
\end{proof}

\begin{lemma}[Closed velocity majorants]
\label{f:lem:qc-closed-velocity-majorants}
For every \(n,m\in\Nzero\),
\begin{align}
 \norm{V_n(0)}_{H^m}&\le\mathfrak D_{n,m},
 \label{f:eq:qc-D-bound}\\
 \norm{V_n}_{L^2(I_T;H^m)}&\le\mathfrak X_{n,m},
 \label{f:eq:qc-X-bound}\\
 \sup_{0\le t\le T}\norm{V_n(t)}_{H^m}
 &\le\mathfrak Z_{n,m}.
 \label{f:eq:qc-Z-bound}
\end{align}
For \(m=-1\), equation~\eqref{f:eq:qc-X-bound} remains valid.  Every
right-hand side is finite and depends on \((u_0,\nu,f,n,m,T)\) through the
finite rectangle-matrix suprema in \eqref{f:eq:qc-B-matrix-identity} and the
corresponding finite prescribed-force coordinates.
\end{lemma}

\begin{proof}
Applying \(\Pp\) to the temporal recurrence gives
\begin{equation}
 V_{n+1}=\nu\Delta V_n
 -\Pp\nabla\!\cdot\sum_{a=0}^{n}\binom na
 (V_a\otimes V_{n-a})+\Pp F_n.
 \label{f:eq:qc-projected-temporal-recurrence}
\end{equation}
Equation~\eqref{f:eq:qc-product-constant} at order \(m+1\), with
\(\Pp F_n\) retained and evaluated at
\(t=0\), proves \eqref{f:eq:qc-D-bound} by induction through
\eqref{f:eq:qc-Drec}.  This estimate, taken pointwise on \([0,T]\),
proves by induction that
\begin{equation}
 \sup_{0\le t\le T}\norm{V_n(t)}_{H^m}
 \le\mathfrak Y_{n,m}.
 \label{f:eq:qc-Y-bound}
\end{equation}
The base case \(m=0\) is the forced estimate \eqref{f:eq:energy-Linf}.  For \(m\ge1\), both
terms in \eqref{f:eq:qc-Y0} are valid: the first is the adjacent trace bound,
and the second follows from
\[
 u(t)=u_0+\int_0^tV_1(s)\dd s
 \quad\text{in }H^m.
\]

The rectangle with middle order \(\mu(m)\), Sobolev nesting, and
\eqref{f:eq:qc-Y-bound} give respectively
\[
 \norm{V_n}_{L^2(I_T;H^m)}
 \le\mathfrak B_{n,\mu(m)}^{1/2},
 \qquad
 \norm{V_n}_{L^2(I_T;H^m)}
 \le T^{1/2}\mathfrak Y_{n,m_+}.
\]
This proves \eqref{f:eq:qc-X-bound} through \eqref{f:eq:qc-Xraw}; the two
special forced-energy cases follow from \eqref{f:eq:energy-Linf} and
\eqref{f:eq:energy-L2H1}.  The adjacent trace lemma gives
\[
 \sup_{0\le t\le T}\norm{V_n(t)}_{H^m}
 \le(1+T^{-1})^{1/2}\mathfrak B_{n,m}^{1/2},
\]
while the Banach-valued fundamental identity and
\eqref{f:eq:qc-D-bound} give
\[
 \sup_{0\le t\le T}\norm{V_n(t)}_{H^m}
 \le\mathfrak D_{n,m}+T^{1/2}\mathfrak X_{n+1,m}.
\]
Taking the minimum of these bounds and \eqref{f:eq:qc-Y-bound} proves
\eqref{f:eq:qc-Z-bound}.
\end{proof}

\begin{definition}[Closed pressure majorants]
\label{f:def:qc-closed-pressure-majorants}
For \(n,m\in\Nzero\), define
\begin{align}
 \mathfrak Q^{(1)}_{n,m}
 &:={}
 A_m\sum_{a=0}^{n}\binom na
 \mathfrak X_{a,m+1}\mathfrak X_{n-a,m+1}
 +\mathfrak L^{(1)}_{n,m},
 \label{f:eq:qc-Q1}\\
 \mathfrak Q^{(2)}_{n,m}
 &:={}
 A_m\sum_{a=0}^{n}\binom na
 \min\left\{
 \mathfrak Z_{a,m+1}\mathfrak X_{n-a,m+1},
 \mathfrak X_{a,m+1}\mathfrak Z_{n-a,m+1}
 \right\}
 +\mathfrak L^{(2)}_{n,m},
 \label{f:eq:qc-Q2}\\
 \mathfrak Q^{(\infty)}_{n,m}
 &:={}
 A_m\sum_{a=0}^{n}\binom na
 \mathfrak Z_{a,m+1}\mathfrak Z_{n-a,m+1}
 +\mathfrak L^{(\infty)}_{n,m},
 \label{f:eq:qc-Qinfty}\\
 \mathfrak Q^{(\mathrm{mod})}_{n,m}
 &:={}
 A_m\sum_{a=0}^{n}\binom na
 \left(
 \mathfrak X_{a+1,m+1}\mathfrak Z_{n-a,m+1}
 +\mathfrak Z_{a,m+1}\mathfrak X_{n-a+1,m+1}
 \right)
 +\mathfrak L^{(2)}_{n+1,m}.
 \label{f:eq:qc-Qmod}
\end{align}
These quantities carry the arguments \([u_0,\nu,f;T]\) through
\(\mathfrak X\) and \(\mathfrak Z\).
\end{definition}

\subsection{Mixed velocity and pressure rows}

The canonical majorants now turn each finite rectangle into bounds for the
actual differentiated velocity and pressure rows, including their traces and
time moduli.

\begin{theorem}[Quantitative mixed-jet rows and traces]
\label{f:thm:qc-mixed-jet}
\begin{samepage}
Let \(N,s\in\Nzero\), let \(0\le n\le N\), and let
\(\alpha\in\Nzero^3\).  Put
\[
 J_{n,\alpha}:=\partial_x^\alpha V_n.
\]
Then
\begin{align}
 \norm{J_{n,\alpha}}_{L^2(I_T;H^{s+1})}^2
 +\norm{J_{n+1,\alpha}}_{L^2(I_T;H^{s-1})}^2
 &\le \mathfrak P_{N,s+|\alpha|},
 \label{f:eq:qc-mixed-adjacent}\\
 \sup_{0\le t\le T}\norm{J_{n,\alpha}(t)}_{H^s}
 &\le \mathfrak T_{N,s+|\alpha|}.
 \label{f:eq:qc-mixed-trace}
\end{align}
\end{samepage}
There is a unique trace \(J_{n,\alpha}^T\in H^s\) at \(t=T\), and
for \(0\le r\le t\le T\),
\begin{align}
 \norm{J_{n,\alpha}(t)-J_{n,\alpha}(r)}_{H^{s-1}}
 &\le (t-r)^{1/2}\mathfrak P_{N,s+|\alpha|}^{1/2},
 \label{f:eq:qc-mixed-low-modulus}\\
 \norm{J_{n,\alpha}(t)-J_{n,\alpha}(r)}_{H^s}
 &\le (t-r)^{1/2}
        \mathfrak P_{N,s+|\alpha|+1}^{1/2},
 \label{f:eq:qc-mixed-strong-modulus}\\
 \norm{J_{n,\alpha}^T}_{H^s}
 &\le \mathfrak T_{N,s+|\alpha|}.
 \label{f:eq:qc-mixed-endpoint-size}
\end{align}
If \(0\le\theta\le1\), then
\begin{equation}
 \norm{J_{n,\alpha}(t)-J_{n,\alpha}(r)}_{H^{s-\theta}}
 \le
 2^{1-\theta}
 \mathfrak T_{N,s+|\alpha|}^{1-\theta}
 \mathfrak P_{N,s+|\alpha|}^{\theta/2}
 (t-r)^{\theta/2}.
 \label{f:eq:qc-mixed-interpolated-modulus}
\end{equation}
For \(1\le a\le\infty\), \(2\le b\le\infty\),
\(k,q\in\Nzero\), and
\(q>k+\frac32-\frac3b\),
\begin{align}
 \norm{\nabla^kJ_{n,\alpha}}_{L^a(I_T;L^b)}
 &\le S_{q,k,b}T^{1/a}\mathfrak T_{N,q+|\alpha|},
 \label{f:eq:qc-mixed-la-lb}\\
 \norm{\nabla^kJ_{n,\alpha}}_{L^2(I_T;L^b)}
 &\le S_{q,k,b}\mathfrak P_{N,q+|\alpha|}^{1/2},
 \label{f:eq:qc-mixed-l2-lb}
\end{align}
where \(T^{1/\infty}=1\).
\end{theorem}

\begin{proof}
Take the rectangle with middle order \(M=s+|\alpha|\).
\begin{samepage}
The Fourier
multiplier inequality
\[
 \norm{\partial_x^\alpha f}_{H^r}
 \le \norm f_{H^{r+|\alpha|}}
\]
and \eqref{f:eq:qc-producer-constant} give
\eqref{f:eq:qc-mixed-adjacent}.
The adjacent-row trace estimate gives
\begin{align*}
 \sup_{0\le t\le T}\norm{J_{n,\alpha}(t)}_{H^s}^2
 &\le T^{-1}\norm{J_{n,\alpha}}_{L^2(I_T;H^s)}^2\\
 &\quad+2\norm{J_{n,\alpha}}_{L^2(I_T;H^{s+1})}
          \norm{J_{n+1,\alpha}}_{L^2(I_T;H^{s-1})}\\
 &\le (T^{-1}+1)\mathfrak P_{N,s+|\alpha|},
\end{align*}
which proves \eqref{f:eq:qc-mixed-trace} and
\eqref{f:eq:qc-mixed-endpoint-size}.
\par
\end{samepage}
The Bochner identity
\[
 J_{n,\alpha}(t)-J_{n,\alpha}(r)
 =\int_r^tJ_{n+1,\alpha}(\tau)\,d\tau
\]
in \(H^{s-1}\), followed by Cauchy--Schwarz, proves
\eqref{f:eq:qc-mixed-low-modulus}.  Applying this argument to the rectangle
with middle order \(s+|\alpha|+1\) proves
\eqref{f:eq:qc-mixed-strong-modulus}.  Interpolation between
\(H^{s-1}\) and \(H^s\), together with
\(\norm{J(t)-J(r)}_{H^s}\le2\mathfrak T_{N,s+|\alpha|}\), proves
\eqref{f:eq:qc-mixed-interpolated-modulus}.  Sobolev embedding and
\eqref{f:eq:qc-mixed-trace} prove \eqref{f:eq:qc-mixed-la-lb}; Sobolev embedding
and \eqref{f:eq:qc-mixed-adjacent} prove \eqref{f:eq:qc-mixed-l2-lb}.
\end{proof}

\begin{theorem}[Quantitative pressure rows]
\label{f:thm:qc-pressure-rows}
\begin{samepage}
Let \(N,s\in\Nzero\), \(0\le n\le N\), and
\(\alpha\in\Nzero^3\).  Set
\[
 \Pi_{n,\alpha}:=\partial_x^\alpha Q_n,
 \qquad r:=s+|\alpha|.
\]
Then
\begin{align}
 \norm{\Pi_{n,\alpha}}_{L^1(I_T;H^s)}
 +\norm{\nabla\Pi_{n,\alpha}}_{L^1(I_T;H^{s-1})}
 &\le 2A_r2^n\mathfrak P_{N,r}+\mathfrak L^{(1)}_{n,r},
 \label{f:eq:qc-pressure-l1}\\
 \norm{\Pi_{n,\alpha}}_{L^2(I_T;H^s)}
 &\le A_r2^n
 \mathfrak T_{N,r+1}\mathfrak P_{N,r}^{1/2}
 +\mathfrak L^{(2)}_{n,r},
 \label{f:eq:qc-pressure-l2}\\
 \sup_{0\le t\le T}\norm{\Pi_{n,\alpha}(t)}_{H^s}
 &\le A_r2^n\mathfrak T_{N,r+1}^2
 +\mathfrak L^{(\infty)}_{n,r}.
 \label{f:eq:qc-pressure-sup}
\end{align}
\end{samepage}
There is a unique endpoint pressure row \(\Pi_{n,\alpha}^T\in H^s\), and
\begin{align}
 \norm{\Pi_{n,\alpha}^T}_{H^s}
 &\le A_r2^n\mathfrak T_{N,r+1}^2
 +\mathfrak L^{(\infty)}_{n,r},
 \label{f:eq:qc-pressure-endpoint}\\
 \norm{\Pi_{n,\alpha}(t)-\Pi_{n,\alpha}(r_0)}_{H^s}
 &\le 2^{n+1}A_r\mathfrak T_{N,r+1}
 \mathfrak P_{N,r+2}^{1/2}|t-r_0|^{1/2}
 +\mathfrak L^{(2)}_{n+1,r}|t-r_0|^{1/2}
 \label{f:eq:qc-pressure-modulus}
\end{align}
for \(0\le r_0\le t\le T\).  If
\(1\le a\le\infty\), \(2\le b\le\infty\),
\(k,q\in\Nzero\), and
\(q>k+\frac32-\frac3b\), then
\begin{equation}
 \norm{\nabla^k\Pi_{n,\alpha}}_{L^a(I_T;L^b)}
 \le
 S_{q,k,b}T^{1/a}A_{q+|\alpha|}2^n
 \mathfrak T_{N,q+|\alpha|+1}^2
 +S_{q,k,b}T^{1/a}\mathfrak L^{(\infty)}_{n,q+|\alpha|}.
 \label{f:eq:qc-pressure-la-lb}
\end{equation}
\end{theorem}

\begin{proof}
The normalized pressure recurrence is
\begin{equation}
 Q_n=R_iR_j\sum_{a=0}^{n}\binom na V_{a,i}V_{n-a,j}-\Phi_n.
 \label{f:eq:qc-pressure-recurrence}
\end{equation}
Equation~\eqref{f:eq:qc-product-constant}, Cauchy--Schwarz in time, and
\(\sum_{a=0}^n\binom na=2^n\) give
\begin{equation}
 \norm{Q_n}_{L^1(I_T;H^r)}
 \le A_r2^n\mathfrak P_{N,r}
 +\norm{\Phi_n}_{L^1(I_T;H^r)}.
 \label{f:eq:qc-pressure-L1-sharp}
\end{equation}
The gradient estimate in \eqref{f:eq:qc-pressure-l1} follows from
\(\norm{\nabla h}_{H^{r-1}}\le\norm h_{H^r}\), with the two prescribed
force-potential norms retained in \(\mathfrak L^{(1)}_{n,r}\).
Using one factor in \(L_t^\infty H_x^{r+1}\) and the other in
\(L_t^2H_x^{r+1}\) proves \eqref{f:eq:qc-pressure-l2}.  Using both factors in
\(L_t^\infty H_x^{r+1}\) proves \eqref{f:eq:qc-pressure-sup} and
\eqref{f:eq:qc-pressure-endpoint}.

Subtract \eqref{f:eq:qc-pressure-recurrence} at times \(t\) and \(r_0\), split
each product difference into two terms, apply
\eqref{f:eq:qc-mixed-strong-modulus} in \(H^{r+1}\), and use
\eqref{f:eq:qc-mixed-trace}.  The two product differences and the binomial sum,
together with
\[
 \Phi_n(t)-\Phi_n(r_0)=\int_{r_0}^t\Phi_{n+1}(\tau)\dd\tau,
\]
give \eqref{f:eq:qc-pressure-modulus}.  Equation~\eqref{f:eq:qc-pressure-la-lb}
follows from \eqref{f:eq:qc-pressure-sup} and Sobolev embedding, including the
force-potential coordinate.
\end{proof}

\begin{proposition}[Matrix form of the closed rectangle estimates]
\label{f:prop:qc-closed-rectangle-matrix}
For nonnegative matrices of identical dimensions, write
\(A\preceq_{\rm ent}B\) when every entry of \(A\) is at most the
corresponding entry of \(B\).  When \(T=T_*<\infty\), set the closed
rectangle matrix by monotone exhaustion and use the three terminal scalar norms
by
\begin{align*}
 \mathbf C_{N,M}(T_*)&:=
 \lim_{S\uparrow T_*}\mathbf C_{N,M}(S),\\
 \mathfrak R_{N,M}(u;T_*)&=
 \lim_{S\uparrow T_*}\mathfrak R_{N,M}(u;S),\\
 \mathfrak Q_{N,M}(p;T_*)&:=
 \lim_{S\uparrow T_*}\mathfrak Q_{N,M}(p;S),\\
 \mathfrak F_{N,M}(f;T_*)&:=
 \lim_{S\uparrow T_*}\mathfrak F_{N,M}(f;S).
\end{align*}
All four limits exist in the extended nonnegative reals by monotonicity.
Theorem~\ref{f:thm:datum-rectangles} and
Proposition~\ref{f:prop:rectangle-pressure-completion} make the scalar limits
for velocity and pressure finite; the prescribed force class makes the force
limit finite, and the estimates below make every matrix entry finite.
Then
\begin{equation}
 \mathbf C_{N,M}(T)
 \preceq_{\rm ent}
 \left(
 \begin{array}{cccc}
  \mathfrak X_{n,M+1}
  &\mathfrak X_{n+1,M-1}
  &\mathfrak Q^{(1)}_{n,M}
  &\mathfrak G^{(1)}_{n,M-1}+\mathfrak L^{(1)}_{n,M}
 \end{array}
 \right)_{n=0}^{N}.
 \label{f:eq:qc-closed-rectangle-matrix}
\end{equation}
Consequently,
\begin{align}
 \mathfrak R_{N,M}(u;T)
 &\le\sum_{n=0}^{N}
 \left(\mathfrak X_{n,M+1}^2
       +\mathfrak X_{n+1,M-1}^2\right),
 \label{f:eq:qc-closed-rectangle-frobenius}\\
 \mathfrak Q_{N,M}(p;T)
 &\le\sum_{n=0}^{N}\mathfrak Q^{(1)}_{n,M}.
 \label{f:eq:qc-closed-pressure-column}\\
 \mathfrak F_{N,M}(f;T)
 &\le\sum_{n=0}^{N}
 \left(\mathfrak G^{(1)}_{n,M-1}+\mathfrak L^{(1)}_{n,M}\right).
 \label{f:eq:qc-closed-force-column}
\end{align}
\end{proposition}

\begin{proof}
The first two columns of \eqref{f:eq:qc-closed-rectangle-matrix} are
\eqref{f:eq:qc-X-bound}.  The pressure recurrence
\eqref{f:eq:qc-pressure-recurrence}, the product estimate
\eqref{f:eq:qc-pressure-product-constant}, and Cauchy--Schwarz in time give
the third column through \eqref{f:eq:qc-Q1}.  The fourth column is the
prescribed-force budget itself, bounded by
\eqref{f:eq:qc-force-row-majorants}--\eqref{f:eq:qc-force-pressure-majorants}.
Take the Frobenius norm of the first two columns and sum the last two columns, using
\eqref{f:eq:rectangle-matrix-norm-identities}.
\end{proof}

\subsection{Frequency and critical-scale consequences}

\begin{proposition}[Velocity and pressure frequency tails]
\label{f:prop:qc-frequency-tails}
Let \(\mathsf P_{\ge K}\) be the Fourier multiplier of
\(\mathbf 1_{\{|\xi|\ge K\}}\), where \(K\ge1\).  Let
\(r_1\in\N\), \(s\in\Nzero\), \(0\le n\le N\), and
\(\alpha\in\Nzero^3\).  Then
\begin{align}
 \norm{\mathsf P_{\ge K}J_{n,\alpha}}_{L^2(I_T;H^s)}
 &\le K^{-r_1}\mathfrak P_{N,s+r_1+|\alpha|}^{1/2},
 \label{f:eq:qc-velocity-tail-l2}\\
 \sup_{0\le t\le T}
 \norm{\mathsf P_{\ge K}J_{n,\alpha}(t)}_{H^s}
 &\le K^{-r_1}\mathfrak T_{N,s+r_1+|\alpha|},
 \label{f:eq:qc-velocity-tail-sup}\\
 \norm{\mathsf P_{\ge K}\Pi_{n,\alpha}}_{L^1(I_T;H^s)}
 &\le K^{-r_1}A_{s+r_1+|\alpha|}2^n
 \mathfrak P_{N,s+r_1+|\alpha|}
 +K^{-r_1}\mathfrak L^{(1)}_{n,s+r_1+|\alpha|},
 \label{f:eq:qc-pressure-tail-l1}\\
 \sup_{0\le t\le T}
 \norm{\mathsf P_{\ge K}\Pi_{n,\alpha}(t)}_{H^s}
 &\le K^{-r_1}A_{s+r_1+|\alpha|}2^n
 \mathfrak T_{N,s+r_1+|\alpha|+1}^2
 +K^{-r_1}\mathfrak L^{(\infty)}_{n,s+r_1+|\alpha|}.
 \label{f:eq:qc-pressure-tail-sup}
\end{align}
The pointwise estimates \eqref{f:eq:qc-velocity-tail-sup} and
\eqref{f:eq:qc-pressure-tail-sup} hold at \(t=T\).
\end{proposition}

\begin{proof}
For every \(f\in H^{s+r_1}\),
\[
 \norm{\mathsf P_{\ge K}f}_{H^s}
 \le K^{-r_1}\norm f_{H^{s+r_1}}.
\]
Apply this inequality to
\eqref{f:eq:qc-mixed-adjacent} and \eqref{f:eq:qc-mixed-trace} at spatial order
\(s+r_1\) to obtain the two velocity bounds.  Since Fourier multipliers
commute,
\[
 \norm{\mathsf P_{\ge K}\Pi_{n,\alpha}}_{H^s}
 \le K^{-r_1}\norm{Q_n}_{H^{s+r_1+|\alpha|}}.
\]
Integration in time followed by
the forced pressure bound \eqref{f:eq:qc-pressure-L1-sharp} at order
\(s+r_1+|\alpha|\) gives
\eqref{f:eq:qc-pressure-tail-l1}.  Applying the tail inequality directly to
\(\Pi_{n,\alpha}\), followed by \eqref{f:eq:qc-pressure-sup} with \(s\) replaced
by \(s+r_1\), gives \eqref{f:eq:qc-pressure-tail-sup}.
\end{proof}

\begin{proposition}[Fixed weighted dyadic shell sums]
\label{f:prop:qc-dyadic-shell-sums}
Let \((\Delta_j)_{j\ge-1}\) be a fixed inhomogeneous
Littlewood--Paley decomposition on \(\R^3\).  Let
\(\sigma\in\R\), \(k\in\Nzero\), and choose \(M\in\Nzero\) such that
\begin{equation}
 M>\sigma+k+\frac32.
 \label{f:eq:qc-dyadic-order-condition}
\end{equation}
Then
\begin{equation}
 \int_0^T\sum_{j\ge-1}2^{\sigma j}
 \norm{\Delta_j\nabla^ku(t)}_{L^\infty}\dd t
 \le C_{\Delta,\sigma,k,M}T^{1/2}\mathfrak X_{0,M}.
 \label{f:eq:qc-dyadic-shell-sum}
\end{equation}
In particular, with \(\omega=\nabla\times u\),
\begin{equation}
 \sum_{j\ge-1}\int_0^T
 \norm{\Delta_j\omega(t)}_{L^\infty}\dd t
 \le C_\Delta T^{1/2}\mathfrak X_{0,3}<\infty.
 \label{f:eq:qc-dyadic-vorticity-sum}
\end{equation}
Each estimate uses one fixed finite Sobolev coordinate; no summation over
Sobolev orders is asserted.
\end{proposition}

\begin{proof}
Bernstein's inequality on a shell gives, for \(j\ge0\),
\[
 2^{\sigma j}\norm{\Delta_j\nabla^kf}_{L^\infty}
 \le C_\Delta 2^{(\sigma+k+3/2)j}\norm{\Delta_jf}_2.
\]
Put \(d:=M-\sigma-k-3/2>0\).  Cauchy--Schwarz in \(j\) yields
\begin{align*}
 \sum_{j\ge0}2^{\sigma j}
 \norm{\Delta_j\nabla^kf}_{L^\infty}
 &\le C_\Delta
 \left(\sum_{j\ge0}2^{-2dj}\right)^{1/2}
 \left(\sum_{j\ge0}2^{2Mj}\norm{\Delta_jf}_2^2\right)^{1/2}\\
 &\le C_{\Delta,\sigma,k,M}\norm f_{H^M}.
\end{align*}
The single low block \(j=-1\) obeys this final bound.  The last
inequality follows from Plancherel's theorem and the finite overlap of the
dyadic multipliers.  Integrating in time, applying Cauchy--Schwarz, and using
\eqref{f:eq:qc-X-bound} prove \eqref{f:eq:qc-dyadic-shell-sum}.  Since curl
commutes with every \(\Delta_j\), the choice
\((\sigma,k,M)=(0,1,3)\) gives
\eqref{f:eq:qc-dyadic-vorticity-sum}.
\end{proof}

\begin{proposition}[Critical-height bounds]
\label{f:prop:qc-critical-height}
Define
\[
 E_\sigma(t):=\frac12\norm{\Lambda^\sigma u(t)}_2^2,
 \qquad
 H(t):=E_{1/2}(t),
\]
and
\[
 E_{j+1/2,\ge K}(t)
 :=\frac12\int_{|\xi|\ge K}|\xi|^{2j+1}
 |\widehat u(t,\xi)|^2\,d\xi.
\]
\begin{samepage}
For \(j\in\Nzero\),
\begin{align}
 \sup_{0\le t\le T}E_{j+1/2}(t)
 &\le\frac12\mathfrak T_{0,j+1}^2,
 \label{f:eq:qc-critical-sup}\\
 \int_0^T E_{j+1/2}(t)\,dt
 &\le\frac12\mathfrak P_{0,j}.
 \label{f:eq:qc-critical-mass}
\end{align}
\end{samepage}
If \(m>j+\frac12\) is an integer and \(K\ge1\), then
\begin{align}
 \sup_{0\le t\le T}E_{j+1/2,\ge K}(t)
 &\le\frac12K^{2j+1-2m}\mathfrak T_{0,m}^2,
 \label{f:eq:qc-critical-tail-sup}\\
 \int_0^T E_{j+1/2,\ge K}(t)\,dt
 &\le\frac12K^{2j+1-2m}\mathfrak P_{0,m-1}.
 \label{f:eq:qc-critical-tail-mass}
\end{align}
\begin{samepage}
For \(k\in\N\),
\begin{equation}
 \norm{H^{(k)}}_{L^1(0,T)}
 \le2^{k-1}\mathfrak P_{k,0}.
 \label{f:eq:qc-height-derivatives}
\end{equation}
In particular,
\begin{align}
 \norm{H'}_{L^1(0,T)}
 &\le\frac12\mathfrak P_{0,2},
 \label{f:eq:qc-critical-variation}\\
 \int_0^T\norm{\Lambda^{3/2}u(t)}_2^2\,dt
 &\le\mathfrak P_{0,1},
 \label{f:eq:qc-critical-dissipation}\\
 \norm{\mathcal P_{1/2}}_{L^1(0,T)}
 &\le\frac12\mathfrak P_{0,2}+\nu\mathfrak P_{0,1}
 +\mathfrak P_{0,0}^{1/2}\mathfrak G^{(2)}_{0,1},
 \label{f:eq:qc-critical-current}\\
 |H(T)-H(t)|
 &\le
 \mathfrak T_{0,1}\mathfrak P_{0,2}^{1/2}(T-t)^{1/2},
 \quad0\le t\le T,
 \label{f:eq:qc-critical-endpoint-modulus}
\end{align}
where
\[
 \mathcal P_{1/2}
 =H'+\nu\norm{\Lambda^{3/2}u}_2^2-\mathcal W_{1/2}.
\]
\par
\end{samepage}
\end{proposition}

\begin{proof}
Equations~\eqref{f:eq:qc-critical-sup} and
\eqref{f:eq:qc-critical-mass} follow from
\(\dot H^{j+1/2}\hookleftarrow H^{j+1}\),
\eqref{f:eq:qc-mixed-trace}, and \eqref{f:eq:qc-mixed-adjacent}.
On \(|\xi|\ge K\),
\[
 |\xi|^{2j+1}\le K^{2j+1-2m}|\xi|^{2m},
\]
which proves \eqref{f:eq:qc-critical-tail-sup} and
\eqref{f:eq:qc-critical-tail-mass}.

The identity
\[
 H^{(k)}(t)
 =\frac12\sum_{a=0}^{k}\binom ka
 \operatorname{Re}\ip{\Lambda^{1/2}V_a(t)}
 {\Lambda^{1/2}V_{k-a}(t)}
\]
and Cauchy--Schwarz in time give
\eqref{f:eq:qc-height-derivatives}.  At \(k=1\), the two equal terms give
\[
 \norm{H'}_{L^1(0,T)}
 \le\norm{u_t}_{L^2(I_T;\dot H^{1/2})}
     \norm u_{L^2(I_T;\dot H^{1/2})}
 \le\frac12\mathfrak P_{0,2}.
\]
Equation~\eqref{f:eq:qc-critical-dissipation} follows from the upper velocity
row of the \((0,1)\) rectangle.  Cauchy--Schwarz gives
\[
 \norm{\mathcal W_{1/2}}_{L^1(0,T)}
 \le\norm u_{L^2(I_T;H^1)}
     \norm{\Pp f}_{L^2(I_T;H^1)}
 \le\mathfrak P_{0,0}^{1/2}\mathfrak G^{(2)}_{0,1}.
\]
The forced height identity defining \(\mathcal P_{1/2}\) then gives
\eqref{f:eq:qc-critical-current}.  Finally,
\[
 \norm{u(T)-u(t)}_{\dot H^{1/2}}
 \le\norm{u(T)-u(t)}_{H^1}
 \le(T-t)^{1/2}\norm{u_t}_{L^2(t,T;H^1)}
 \le(T-t)^{1/2}\mathfrak P_{0,2}^{1/2},
\]
and
\[
 |H(T)-H(t)|
 \le\frac12\bigl(
 \norm{u(T)}_{\dot H^{1/2}}+\norm{u(t)}_{\dot H^{1/2}}
 \bigr)\norm{u(T)-u(t)}_{\dot H^{1/2}}.
\]
Equation~\eqref{f:eq:qc-mixed-trace} bounds each retained velocity factor, so
these two estimates prove \eqref{f:eq:qc-critical-endpoint-modulus}.
\end{proof}

Taking \(j=0\) and any integer \(m\ge1\) in
\eqref{f:eq:qc-critical-tail-sup} gives the uniform tail limit
\begin{equation}
 \lim_{K\to\infty}\sup_{0\le t\le T}
 E_{1/2,\ge K}(t)=0.
 \label{f:eq:qc-critical-tail-extinction}
\end{equation}

\subsection{Continuation norms}

\begin{proposition}[Continuation norms]
\label{f:prop:qc-continuation-norms}
The finite-horizon solution satisfies
\begin{align}
 \norm u_{L^2(I_T;L^\infty)}
 &\le S_{2,0,\infty}\mathfrak P_{0,1}^{1/2},
 \label{f:eq:qc-l2-linfty}\\
 \int_0^T\norm{\nabla u(t)}_{L^\infty}\,dt
 &\le S_{3,1,\infty}T^{1/2}\mathfrak P_{0,2}^{1/2},
 \label{f:eq:qc-lipschitz-action}\\
 \int_0^T\norm{\nabla\times u(t)}_{L^\infty}\,dt
 &\le 2S_{3,1,\infty}T^{1/2}\mathfrak P_{0,2}^{1/2}.
 \label{f:eq:qc-vorticity-action}
\end{align}
For every integer \(m\ge3\), let \(K_m\) be the explicit constant in
\eqref{f:eq:explicit-commutator-constant}.  Then
\begin{equation}
 \left|\operatorname{Re}
 \ip{A^m\Pp\nabla\!\cdot(u\otimes u)}{A^m u}_{L^2}\right|
 \le K_m\norm{\nabla u}_{L^\infty}\norm u_{H^m}^2.
 \label{f:eq:qc-commutator-constant}
\end{equation}
\begin{equation}
 \sup_{0\le t\le T}\norm{u(t)}_{H^m}^2
 \le
 \left(\norm{u_0}_{H^m}+\mathfrak G^{(1)}_{0,m}\right)^2
 \exp\!\left(
 2K_mS_{3,1,\infty}T^{1/2}\mathfrak P_{0,2}^{1/2}
 \right).
 \label{f:eq:qc-continuation-growth}
\end{equation}
\end{proposition}

\begin{proof}
Sobolev embedding and the upper rows of the \((0,1)\) and \((0,2)\)
rectangles give \eqref{f:eq:qc-l2-linfty} and
\eqref{f:eq:qc-lipschitz-action}; the pointwise bound
\(\norm{\nabla\times u}_\infty\le2\norm{\nabla u}_\infty\) gives
\eqref{f:eq:qc-vorticity-action}.  The forced \(H^m\) energy identity and
\eqref{f:eq:qc-commutator-constant} give
\[
 \frac12\frac d{dt}\norm u_{H^m}^2
 +\nu\norm{\nabla u}_{H^m}^2
 \le K_m\norm{\nabla u}_{L^\infty}\norm u_{H^m}^2
 +\norm{\Pp f}_{H^m}\norm u_{H^m}.
\]
Regularizing the norm by
\((\norm u_{H^m}^2+\varepsilon)^{1/2}\), applying Gronwall, and then
letting \(\varepsilon\downarrow0\) gives
\[
 \norm{u(t)}_{H^m}
 \le\left(\norm{u_0}_{H^m}
 +\int_0^t\norm{\Pp f(s)}_{H^m}\dd s\right)
 \exp\!\left(K_m\int_0^t\norm{\nabla u(s)}_\infty\dd s\right).
\]
Gronwall's inequality and \eqref{f:eq:qc-lipschitz-action} prove
\eqref{f:eq:qc-continuation-growth}.
\end{proof}

\newcommand{\ForcedQuantitativeCatalogue}{%
\section{Forced canonical coordinate catalogue}
\label{f:sec:quantitative-catalogue}

The estimates in Section~\ref{f:sec:quantitative-consequences} are organized by
their role in the argument.  For reference, the following theorem collects
their canonical-coordinate forms in one place.

\begin{theorem}[Canonical finite-horizon coordinate estimates]
\label{f:thm:qc-sharp-datum-estimates}
Under the hypotheses and notation of
Definition~\ref{f:def:qc-producer-constants}, let
\(J_{n,\alpha}:=\partial_x^\alpha V_n\) and
\(\Pi_{n,\alpha}:=\partial_x^\alpha Q_n\), as in
Theorems~\ref{f:thm:qc-mixed-jet} and \ref{f:thm:qc-pressure-rows}.
Let \(n,s,k,q\in\Nzero\), \(\alpha\in\Nzero^3\), and set
\(r:=s+|\alpha|\).  Then
\begin{align}
 \norm{J_{n,\alpha}}_{L^2(I_T;H^{s+1})}
 &\le\mathfrak X_{n,r+1},
 \label{f:eq:qc-sharp-mixed-upper}\\
 \norm{J_{n+1,\alpha}}_{L^2(I_T;H^{s-1})}
 &\le\mathfrak X_{n+1,r-1},
 \label{f:eq:qc-sharp-mixed-lower}\\
 \sup_{0\le t\le T}\norm{J_{n,\alpha}(t)}_{H^s}
 &\le\mathfrak Z_{n,r},
 \label{f:eq:qc-sharp-mixed-sup}\\
 \norm{J_{n,\alpha}^T}_{H^s}
 &\le\mathfrak Z_{n,r},
 \label{f:eq:qc-sharp-mixed-endpoint}
\end{align}
\begin{samepage}
\begin{align}
 \norm{J_{n,\alpha}(t)-J_{n,\alpha}(r_0)}_{H^{s-1}}
 &\le\min\left\{
 |t-r_0|^{1/2}\mathfrak X_{n+1,r-1},
 |t-r_0|\mathfrak Z_{n+1,r}
 \right\},
 \label{f:eq:qc-sharp-mixed-low-modulus}\\
 \norm{J_{n,\alpha}(t)-J_{n,\alpha}(r_0)}_{H^s}
 &\le\min\left\{
 |t-r_0|^{1/2}\mathfrak X_{n+1,r},
 |t-r_0|\mathfrak Z_{n+1,r}
 \right\},
 \label{f:eq:qc-sharp-mixed-modulus}
\end{align}
for \(0\le r_0\le t\le T\).
\par
\end{samepage}
For \(0\le\theta\le1\),
\begin{equation}
 \norm{J_{n,\alpha}(t)-J_{n,\alpha}(r_0)}_{H^{s-\theta}}
 \le2^{1-\theta}\mathfrak Z_{n,r}^{1-\theta}
 \mathfrak X_{n+1,r-1}^{\theta}|t-r_0|^{\theta/2}.
 \label{f:eq:qc-sharp-mixed-interpolation}
\end{equation}
\begin{samepage}
If \(1\le a\le\infty\), \(2\le b\le\infty\), and
\(q>k+\frac32-\frac3b\), then
\begin{align}
 \norm{\nabla^kJ_{n,\alpha}}_{L^a(I_T;L^b)}
 &\le S_{q,k,b}T^{1/a}\mathfrak Z_{n,q+|\alpha|},
 \label{f:eq:qc-sharp-mixed-Lab}\\
 \norm{\nabla^kJ_{n,\alpha}}_{L^2(I_T;L^b)}
 &\le S_{q,k,b}\mathfrak X_{n,q+|\alpha|},
 \label{f:eq:qc-sharp-mixed-L2b}
\end{align}
where \(T^{1/\infty}=1\).
\end{samepage}

The pressure rows obey
\begin{align}
 \norm{\Pi_{n,\alpha}}_{L^1(I_T;H^s)}
 +\norm{\nabla\Pi_{n,\alpha}}_{L^1(I_T;H^{s-1})}
 &\le\mathfrak Q^{(1)}_{n,r},
 \label{f:eq:qc-sharp-pressure-L1}\\
 \norm{\Pi_{n,\alpha}}_{L^2(I_T;H^s)}
 &\le\mathfrak Q^{(2)}_{n,r},
 \label{f:eq:qc-sharp-pressure-L2}\\
 \sup_{0\le t\le T}\norm{\Pi_{n,\alpha}(t)}_{H^s}
 &\le\mathfrak Q^{(\infty)}_{n,r},
 \label{f:eq:qc-sharp-pressure-sup}\\
 \norm{\Pi_{n,\alpha}^T}_{H^s}
 &\le\mathfrak Q^{(\infty)}_{n,r},
 \label{f:eq:qc-sharp-pressure-endpoint}\\
 \norm{\Pi_{n,\alpha}(t)-\Pi_{n,\alpha}(r_0)}_{H^s}
 &\le\min\left\{
 |t-r_0|^{1/2}\mathfrak Q^{(\mathrm{mod})}_{n,r},
 |t-r_0|\mathfrak Q^{(\infty)}_{n+1,r}
 \right\}.
 \label{f:eq:qc-sharp-pressure-modulus}
\end{align}
Under the preceding restrictions on \(a,b,k,q\),
\begin{equation}
 \norm{\nabla^k\Pi_{n,\alpha}}_{L^a(I_T;L^b)}
 \le S_{q,k,b}T^{1/a}
 \mathfrak Q^{(\infty)}_{n,q+|\alpha|}.
 \label{f:eq:qc-sharp-pressure-Lab}
\end{equation}

Let \(r_1\in\N\) and \(K\ge1\).  Then
\begin{align}
 \norm{\mathsf P_{\ge K}J_{n,\alpha}}_{L^2(I_T;H^s)}
 &\le K^{-r_1}\mathfrak X_{n,s+r_1+|\alpha|},
 \label{f:eq:qc-sharp-velocity-tail-L2}\\
 \sup_{0\le t\le T}
 \norm{\mathsf P_{\ge K}J_{n,\alpha}(t)}_{H^s}
 &\le K^{-r_1}\mathfrak Z_{n,s+r_1+|\alpha|},
 \label{f:eq:qc-sharp-velocity-tail-sup}\\
 \norm{\mathsf P_{\ge K}\Pi_{n,\alpha}}_{L^1(I_T;H^s)}
 &\le K^{-r_1}\mathfrak Q^{(1)}_{n,s+r_1+|\alpha|},
 \label{f:eq:qc-sharp-pressure-tail-L1}\\
 \sup_{0\le t\le T}
 \norm{\mathsf P_{\ge K}\Pi_{n,\alpha}(t)}_{H^s}
 &\le K^{-r_1}\mathfrak Q^{(\infty)}_{n,s+r_1+|\alpha|}.
 \label{f:eq:qc-sharp-pressure-tail-sup}
\end{align}

For \(j\in\Nzero\),
\begin{align}
 \sup_{0\le t\le T}E_{j+1/2}(t)
 &\le\frac12\mathfrak Z_{0,j+1}^2,
 \label{f:eq:qc-sharp-critical-sup}\\
 \int_0^TE_{j+1/2}(t)\dd t
 &\le\frac12\mathfrak X_{0,j+1}^2.
 \label{f:eq:qc-sharp-critical-mass}
\end{align}
If \(m>j+\frac12\) is an integer, then
\begin{align}
 \sup_{0\le t\le T}E_{j+1/2,\ge K}(t)
 &\le\frac12K^{2j+1-2m}\mathfrak Z_{0,m}^2,
 \label{f:eq:qc-sharp-critical-tail-sup}\\
 \int_0^TE_{j+1/2,\ge K}(t)\dd t
 &\le\frac12K^{2j+1-2m}\mathfrak X_{0,m}^2.
 \label{f:eq:qc-sharp-critical-tail-mass}
\end{align}
For \(k\in\N\),
\begin{equation}
 \norm{H^{(k)}}_{L^1(0,T)}
 \le\frac12\sum_{a=0}^{k}\binom ka
 \mathfrak X_{a,1}\mathfrak X_{k-a,1}.
 \label{f:eq:qc-sharp-height-derivatives}
\end{equation}
\begin{samepage}
For \(j,\ell\in\Nzero\), define
\begin{align}
 \mathfrak E_{j,\ell}^{[1]}
 &:=\frac12\sum_{a=0}^{\ell}\binom\ell a
 \mathfrak X_{a,j+1}\mathfrak X_{\ell-a,j+1},
 \notag\\
 \mathfrak E_{j,\ell}^{[\infty]}
 &:=\frac12\sum_{a=0}^{\ell}\binom\ell a
 \mathfrak Z_{a,j+1}\mathfrak Z_{\ell-a,j+1},
 \label{f:eq:qc-height-row-majorants}\\
 \mathfrak W_{j,\ell}^{[1]}
 &:=\sum_{a=0}^{\ell}\binom\ell a
 \mathfrak X_{a,j+1}\mathfrak G^{(2)}_{\ell-a,j+1},
 \notag\\
 \mathfrak W_{j,\ell}^{[\infty]}
 &:=\sum_{a=0}^{\ell}\binom\ell a
 \mathfrak Z_{a,j+1}\mathfrak G^{(\infty)}_{\ell-a,j+1},
 \label{f:eq:qc-force-work-row-majorants}\\
 \mathfrak C_{j,\ell}^{[1]}
 &:=\mathfrak E_{j,\ell+1}^{[1]}
    +2\nu\mathfrak E_{j+1,\ell}^{[1]}
    +\mathfrak W_{j,\ell}^{[1]},
 \notag\\
 \mathfrak C_{j,\ell}^{[\infty]}
 &:=\mathfrak E_{j,\ell+1}^{[\infty]}
    +2\nu\mathfrak E_{j+1,\ell}^{[\infty]}
    +\mathfrak W_{j,\ell}^{[\infty]}.
 \label{f:eq:qc-production-row-majorants}
\end{align}
Then the complete differentiated energy and production rows obey
\begin{align}
 \norm{\partial_t^\ell E_{j+1/2}}_{L^1(0,T)}
 &\le\mathfrak E_{j,\ell}^{[1]},
 &
 \sup_{0\le t\le T}
 \abs{\partial_t^\ell E_{j+1/2}(t)}
 &\le\mathfrak E_{j,\ell}^{[\infty]},
 \label{f:eq:qc-sharp-energy-row-family}\\
 \norm{\partial_t^\ell\mathcal W_{j+1/2}}_{L^1(0,T)}
 &\le\mathfrak W_{j,\ell}^{[1]},
 &
 \sup_{0\le t\le T}
 \abs{\partial_t^\ell\mathcal W_{j+1/2}(t)}
 &\le\mathfrak W_{j,\ell}^{[\infty]},
 \label{f:eq:qc-sharp-force-work-row-family}\\
 \norm{\partial_t^\ell\cP_{j+1/2}}_{L^1(0,T)}
 &\le\mathfrak C_{j,\ell}^{[1]},
 &
 \sup_{0\le t\le T}
 \abs{\partial_t^\ell\cP_{j+1/2}(t)}
 &\le\mathfrak C_{j,\ell}^{[\infty]}.
 \label{f:eq:qc-sharp-production-row-family}
\end{align}
\end{samepage}
Every current in \eqref{f:eq:completed-height-recurrence} now has a
finite-order majorant retaining temporal depth, spatial height, viscosity,
and horizon.
In particular,
\begin{align}
 \norm{H'}_{L^1(0,T)}
 &\le\mathfrak X_{1,1}\mathfrak X_{0,1},
 \label{f:eq:qc-sharp-height-variation}\\
 \int_0^T\norm{\Lambda^{3/2}u(t)}_2^2\dd t
 &\le\mathfrak X_{0,2}^2,
 \label{f:eq:qc-sharp-critical-dissipation}\\
 \norm{\mathcal P_{1/2}}_{L^1(0,T)}
 &\le\mathfrak X_{1,1}\mathfrak X_{0,1}
 +\nu\mathfrak X_{0,2}^2
 +\mathfrak X_{0,1}\mathfrak G^{(2)}_{0,1},
 \label{f:eq:qc-sharp-critical-current}\\
 |H(T)-H(t)|
 &\le\min\left\{
 \mathfrak Z_{0,1}\mathfrak X_{1,1}(T-t)^{1/2},
 \mathfrak Z_{0,1}\mathfrak Z_{1,1}(T-t)
 \right\},
 \qquad 0\le t\le T.
 \label{f:eq:qc-sharp-critical-modulus}
\end{align}

Finally,
\begin{align}
 \norm u_{L^2(I_T;L^\infty)}
 &\le S_{2,0,\infty}\mathfrak X_{0,2},
 \label{f:eq:qc-sharp-L2Linfty}\\
 \int_0^T\norm{\nabla u(t)}_\infty\dd t
 &\le S_{3,1,\infty}T^{1/2}\mathfrak X_{0,3},
 \label{f:eq:qc-sharp-lipschitz-action}\\
 \int_0^T\norm{\nabla\times u(t)}_\infty\dd t
 &\le2S_{3,1,\infty}T^{1/2}\mathfrak X_{0,3},
 \label{f:eq:qc-sharp-vorticity-action}\\
 \sup_{0\le t\le T}\norm{u(t)}_{H^m}^2
 &\le\left(\norm{u_0}_{H^m}+\mathfrak G^{(1)}_{0,m}\right)^2
 \exp\!\left(
 2K_mS_{3,1,\infty}T^{1/2}\mathfrak X_{0,3}
 \right),
 \qquad m\ge3.
 \label{f:eq:qc-sharp-continuation-growth}
\end{align}
Every right-hand side is an explicit finite expression in the displayed
analytic constants, the prescribed-force coordinates, the initial-jet recursion
\eqref{f:eq:qc-D0}--\eqref{f:eq:qc-Drec}, and the finite collection of canonical
rectangle-norm majorants \(\mathfrak B_{a,b}[u_0,\nu,f;T]\) selected by the stated
indices.  Equations \eqref{f:eq:qc-producer-through-B} and
\eqref{f:eq:qc-Y00}--\eqref{f:eq:qc-Qmod} specify every intermediate operation.
Equation \eqref{f:eq:qc-B-by-height-reconstruction} bounds each such rectangle
norm through the finite height recursion
\eqref{f:eq:qc-height-producer-coordinate}--%
\eqref{f:eq:qc-height-reconstructed-rectangle}.  Together with the reverse
bound in \eqref{f:eq:qc-height-spatial-bounds}, this identifies the mixed-row
and height descriptions of the original forced history at every finite order.
\begin{samepage}
This is the exact coordinate-level catalogue for general data: every finite
quantity is expressed through the compatible completed-height and mixed-row
coordinates selected by its displayed indices.\par
\end{samepage}
Proposition~\ref{f:prop:qc-exact-viscosity-scaling} displays the exact powers of
\(\nu\) and the rescaled horizon \(\nu T\).  For every fixed choice of the
displayed indices, the definitions select only finitely many rectangle-norm
majorants.
Under the explicit critical smallness condition
\eqref{f:eq:app-critical-smallness}, Proposition~\ref{f:prop:app-small-data-rectangles}
eliminates those majorants and gives finite formulas solely in the
displayed velocity-datum and prescribed-force norms, \(\nu\), \(T\),
\(N\), and \(M\).
\end{theorem}

\begin{proof}
The Fourier multiplier inequality
\[
 \norm{\partial_x^\alpha f}_{H^\ell}
 \le\norm f_{H^{\ell+|\alpha|}},
\]
followed by Lemma~\ref{f:lem:qc-closed-velocity-majorants}, proves
\eqref{f:eq:qc-sharp-mixed-upper}--\eqref{f:eq:qc-sharp-mixed-endpoint}.
The Banach identity
\[
 J_{n,\alpha}(t)-J_{n,\alpha}(r_0)
 =\int_{r_0}^tJ_{n+1,\alpha}(\tau)\dd\tau
\]
in the two indicated Sobolev spaces and Cauchy--Schwarz prove the
square-root branches of \eqref{f:eq:qc-sharp-mixed-low-modulus} and
\eqref{f:eq:qc-sharp-mixed-modulus}.
\begin{samepage}
This identity, followed by the
pointwise bound \eqref{f:eq:qc-Z-bound} for \(V_{n+1}\), proves their linear
branches.  Interpolation between the low modulus
and
\[
 \norm{J_{n,\alpha}(t)-J_{n,\alpha}(r_0)}_{H^s}
 \le2\mathfrak Z_{n,r}
\]
proves \eqref{f:eq:qc-sharp-mixed-interpolation}.
\end{samepage}
The fixed Sobolev embedding constant in \eqref{f:eq:qc-embedding-constant} proves
\eqref{f:eq:qc-sharp-mixed-Lab} and \eqref{f:eq:qc-sharp-mixed-L2b}.

Apply the recurrence
\[
 Q_n=R_iR_j\sum_{a=0}^{n}\binom naV_{a,i}V_{n-a,j}
 -\Phi_n
\]
term by term at order \(r\), using
\eqref{f:eq:qc-pressure-product-constant}.  Cauchy--Schwarz in time gives
\eqref{f:eq:qc-sharp-pressure-L1}; placing either factor in
\(L_t^\infty H_x^{r+1}\) gives
\eqref{f:eq:qc-sharp-pressure-L2}; placing both factors there gives
\eqref{f:eq:qc-sharp-pressure-sup} and
\eqref{f:eq:qc-sharp-pressure-endpoint}, with the corresponding
\(\mathfrak L\) coordinate retaining \(-\Phi_n\) in each case.  After subtracting the recurrence at
times \(t\) and \(r_0\), split each product difference into two terms and
apply \eqref{f:eq:qc-sharp-mixed-modulus}; this gives the square-root branch
of \eqref{f:eq:qc-sharp-pressure-modulus}.  Since
\(\partial_t\Pi_{n,\alpha}=\Pi_{n+1,\alpha}\), the Banach identity and
\eqref{f:eq:qc-sharp-pressure-sup} at row \(n+1\) give its linear branch.
Sobolev embedding proves
\eqref{f:eq:qc-sharp-pressure-Lab}.

For every \(f\in H^{\ell+r_1}\),
\[
 \norm{\mathsf P_{\ge K}f}_{H^\ell}
 \le K^{-r_1}\norm f_{H^{\ell+r_1}}.
\]
Apply this inequality at spatial order \(s+r_1\) to
\eqref{f:eq:qc-X-bound}, \eqref{f:eq:qc-Z-bound},
\eqref{f:eq:qc-sharp-pressure-L1}, and
\eqref{f:eq:qc-sharp-pressure-sup}.  The four resulting estimates are
\eqref{f:eq:qc-sharp-velocity-tail-L2}--%
\eqref{f:eq:qc-sharp-pressure-tail-sup}.

The inequalities
\[
 \norm f_{\dot H^{j+1/2}}\le\norm f_{H^{j+1}},
 \qquad
 |\xi|^{2j+1}\le K^{2j+1-2m}|\xi|^{2m}
 \quad(|\xi|\ge K),
\]
give \eqref{f:eq:qc-sharp-critical-sup}--%
\eqref{f:eq:qc-sharp-critical-tail-mass}.  The exact identity
\[
 H^{(k)}(t)=\frac12\sum_{a=0}^{k}\binom ka
 \operatorname{Re}\ip{\Lambda^{1/2}V_a(t)}
 {\Lambda^{1/2}V_{k-a}(t)}_{L^2}
\]
and Cauchy--Schwarz in time give
\eqref{f:eq:qc-sharp-height-derivatives} and
\eqref{f:eq:qc-sharp-height-variation}.  More generally,
\[
 \partial_t^\ell E_{j+1/2}
 =\frac12\sum_{a=0}^{\ell}\binom\ell a
 \operatorname{Re}\ip{\Lambda^{j+1/2}V_a}
 {\Lambda^{j+1/2}V_{\ell-a}}_{L^2}.
\]
Cauchy--Schwarz in time and the pointwise estimates
\eqref{f:eq:qc-X-bound}--\eqref{f:eq:qc-Z-bound} prove
\eqref{f:eq:qc-sharp-energy-row-family}.  Leibniz' rule applied to
\(\mathcal W_{j+1/2}\) gives the majorants
\eqref{f:eq:qc-force-work-row-majorants}.  Differentiating the forced identity
\[
 \cP_{j+1/2}=\partial_tE_{j+1/2}+2\nu E_{j+3/2}
 -\mathcal W_{j+1/2}
\]
then proves \eqref{f:eq:qc-sharp-production-row-family}.  The next two
estimates follow from
\[
 \mathcal P_{1/2}=H'+\nu\norm{\Lambda^{3/2}u}_2^2
 -\mathcal W_{1/2}.
\]
\begin{samepage}
For the terminal critical-height modulus,
\[
 |H(T)-H(t)|
 \le\frac12\bigl(
 \norm{u(T)}_{\dot H^{1/2}}+\norm{u(t)}_{\dot H^{1/2}}
 \bigr)\norm{u(T)-u(t)}_{\dot H^{1/2}},
\]
and the Banach identity in \(H^1\) prove the square-root branch of
\eqref{f:eq:qc-sharp-critical-modulus}.
\par
\end{samepage}
The identity
\[
 H'=\operatorname{Re}
 \ip{\Lambda^{1/2}V_1}{\Lambda^{1/2}V_0}_{L^2},
\]
together with \eqref{f:eq:qc-Z-bound}, proves its linear branch.

Sobolev embedding and \eqref{f:eq:qc-X-bound} prove
\eqref{f:eq:qc-sharp-L2Linfty}--\eqref{f:eq:qc-sharp-vorticity-action}.
Lemma~\ref{f:lem:transport-commutator}, with the prescribed force retained,
gives after regularization of the norm
\[
 \frac d{dt}\norm u_{H^m}
 \le K_m\norm{\nabla u}_\infty\norm u_{H^m}
 +\norm{\Pp f}_{H^m}.
\]
The inhomogeneous Gronwall inequality and \eqref{f:eq:qc-sharp-lipschitz-action} prove
\eqref{f:eq:qc-sharp-continuation-growth}.
\end{proof}
}
